\emph{By Steve Leigh}

These notes were compiled by a member of Firebrand for the Introduction to Revolutionary Socialism study group:

\paragraph{What Is Imperialism?}

Imperialism is the stage of capitalism in which a few economically advanced states use their economic and military power to subjugate and exploit weaker nations. One of the consequences of imperialism is liberation movements and thus, nationalism. However, we must remember that class struggle, not national identity, is the motor of history. One of the ramifications of nationalism is its divisive; workers in dominant nations come to think they stand to gain from oppressing weaker nations. This article will present Lenin's argument in favor of the right of nations to self-determination.

\paragraph{Marx and Engels on National Oppression}

Initially, Marx and Engels supported many national movements, believing capitalism to be a historic step forward from feudalism. The transition to nation-states would certainly see flaws and exploitation, but at the same time would create a class of producers, the proletarians, with the potential for democratic self-rule. They believed the working class would grow and eventually could democratically govern themselves. Furthermore, differences between nations would decrease as capitalism drove globalization, and socialist revolution within nations could ultimately end the oppression of workers.

These ideas led to their support of national movements (Poland, Hungary, Germany, Italy) against the feudal oppressors of the time (Russia and Austria). Accordingly, they opposed movements of Serbs, Croatians, and Czechs during the 1848 revolution as they were allies of Russia and Austria. Marx and Engels also supported the Indian Revolution (Great Rebellion of 1857) despite believing British imperialism was helping destroy the pre-capitalist mode of production in India; they supported national movements based on a political assessment rather than an economic one.

The Irish national struggle of the 1860s led Marx and Engels to modify their views on national movements, recognizing imperialism as a relationship between oppressor and oppressed nations. They underestimated the ability of less economically developed nations to achieve bourgeois democratic freedoms short of socialist revolution and overlooked the hold of nationalism over workers in more developed nations. They realized that now national movements were no longer about replacing feudalism but about freeing oppressed nations from imperialist powers. Marx believed that the English working class should side with the Irish rather than their own ruling class, not out of sympathy, but for their own material interest. This realization led to Marx and Engels' key insight:

\blockquote{The nationalism of the workers belonging to an oppressor nation binds them to their rulers and only does harm to themselves, while the nationalism of an oppressed nation can lead them to fight back against those rulers.}

\paragraph{The 1903 Congress}

Lenin belonged to the Russian Social Democratic Labor Party (RSDLP) which convened its Second Congress in the summer of 1903. This Congress was an important moment for Marxism for two splits:

The Jewish Labor Bund was a socialist party of the Jewish proletarians who rejected Zionism, recognizing it as a form of reactionary Jewish nationalism. They simultaneously advocated for horrible struggles and oppression of Jewish workers in Russia. The Bund wanted a federation to give them complete autonomy over their interests with no veto/overruling power by the larger RSDLP party. Lenin and Trotsky were against the idea of a federation of solely Jewish workers, fearing it would foster distrust and result in segregation and alienation. He suggested the Party should be organized by shared interest, not by ethnic or other group. He believed that it was the Central Party's responsibility to care about every single member of the party, not just the majorities, and to fight against all forms of oppression. This was the best way to practice \emph{democratic centralism}, where there is the opportunity for internal debate across all groups, but the party acts as one after a majority vote.

The 1903 Congress rejected the Bund's demand to be the sole representatives of the Jewish proletarians and reaffirmed that the RSDLP should unite workers of all nationalities/organizations without exception. The Bund seceded from the party (though many Jewish workers remained active in RSDLP branches).

The second debate concerned the significance of the right of nations to self-determination. Rosa Luxemburg, a member of the Polish Social Democrats, felt that the demand for Polish independence was obsolete because the Russian and Polish economies were so integrated. She instead believed in the freedom of cultural development of each nationality through the democratization of existing states. Lenin disagreed with her. He thought socialism could not be based on economics alone, and that capitalism and imperialism must also be fought politically. He thought allowing a nation to unite against its oppressor and free itself politically was foundational to a successful socialist revolution.

Lenin believed the right of nations is essential; for nations to be able to secede from their oppressors. He envisioned that a state's territories would be self-governed and autonomous. Institutions would be organized through universal, equal, and secret votes of a territory's population, including the decision to secede from the state. Lenin also made clear that socialists should demand unconditional equality for all nationalities, without giving privileges for some nationalities, such as through the establishment of an official state language.

\paragraph{Assessing National Movements and Against Cultural-National Autonomy}

Lenin sees two periods of capitalism. Stage One was generally progressive, as it saw feudalism overthrown and drew all classes of the population into politics. Stage Two saw the development of long-established constitutional regimes and highly developed class divisions/antagonism. Even though socialists must unconditionally support the right of oppressed nations to self-determination, we should always evaluate on a case-by-case basis whether the actual independence of a nation is aligned with the interests of the international working class.

Lenin was against Otto Bauer's idea of cultural-national autonomy, which recognizes an individual cultural basis of autonomy rather than one based on the actual territory within a state. He warned that grouping people culturally would deepen class-based divisions, pointing to ``separate but equal'' Jim Crow laws in the US. Lenin thought we should only take the democratic and socialist elements of a culture, to oppose bourgeois culture and build an international higher culture of unity.

\paragraph{Luxemburg and Imperialist Economism}

Luxemburg thought the actual possibility of self-determination was a utopia. Under capitalism, it wasn't economically possible for smaller nations to be truly independent from the powerful imperialist nations. The development of imperialism undermines the possibility of self-determination because global capitalism will still slowly destroy primitive societies. It was, therefore, hopeless and reactionary to divide all existing states into national units while they still depended on capitalist development. Luxembourg thought the right to self-determination could only be realized within a socialist regime.

Lenin disagreed; as imperialism tightened the grip of stronger nations on weaker ones, so would it heighten the legitimacy of the socialist demand for the right to self-determination. It would be impossible to abolish national oppression under capitalism, since this requires the abolition of classes, but socialism cannot be reduced to economics alone. By transforming capitalism into socialism through revolution, the proletariat creates the possibility of abolishing national oppression and establishing full democracy in all spheres. This would include the delineation of state frontiers according to the desire of the population, including complete freedom to secede. In his debate with Luxemburg on the national question, Lenin emphasized at all times that the

\blockquote{self-determination of nations today hinges on the conduct of socialists in the oppressor nations. A socialist of any of the oppressor nations \ldots who does not recognize and does not struggle for the right of oppressed nations to self-determination (i.e., the right to secession) is in reality a chauvinist, not a socialist.}

\paragraph{Poland and the First World War}

With the First World War looming, Karl Kautsky, a prominent socialist theoretician, failed the test of chauvinism, arguing that every nation, even an imperialist nation, had the right to defend itself.(This is similar to those who support Israel's right to defend itself against the colonized and oppressed Palestinians today) He believed that the carving up of the world by global powers would lead to stability. Lenin was scathing in his attack on socialists in support of the imperialist war; the course of the war proving Lenin right.

In 1915, Germany captured Poland from Russia. Rosa Luxemburg and the Polish Social Democrats argued it did not matter whether Poland was inevitably dominated by Russia or Germany. However, Lenin believed that Russia had benefitted by ceasing to oppress and hold Poland by force, and the German proletariat had lost in terms of internationalism. He recognized that achieving Polish independence from both Germany and Russia was impossible at the moment, short of an all-European war. Thus, he advocated for socialists in oppressor nations to uphold Poland's right to self-determination, even while withholding support from the imperialist power that claimed to be supporting Poland. Meanwhile, Polish socialists should work for the unity of the working class and lay the ground for future independence.

Even when the slogan of independence cannot be raised because of the alignment of a given national movement with one among a set of rival imperialist powers, revolutionary socialists in both the oppressor and oppressed countries must find a way to affirm the right of oppressed nations to self-determination.

\paragraph{Conclusion}

Lenin believed socialists must support a nation's right to self-determination in order to advance a united, international class struggle. He clarified the distinction between the nationalism of the oppressor and the nationalism of the oppressed, and he pointed out that the act of carrying out a liberation struggle, even when imperfect, would form the working-class leaders of the future. He also recognized there would be concessions made and patience required; forming a voluntary union requires trust between workers unlearning capitalist individualism. Lenin's critics argue his idea is contradictory, supporting the `right' to separation without actually advocating its immediate exercise. But Lenin responded that there could not be any other road to internationalism and the amalgamation of nations from the present situation. The right of nations to self-determination is one of the most important achievements of Marxism and a beacon of workers' democracy that guides socialists in anti-imperial struggles.
